% --- Archon physics formalization source begin ---
% archon:physics
% archon:covers PhyXMiniProblems/problem_phyx_mini_0257.lean
% archon:source-report reports/phyx_mini/problem_phyx_mini_0257.source.json

\chapter{Physics problem phyx\_mini\_0257}
\label{ch:PhyXMiniProblems_problem_phyx_mini_0257}

\paragraph{Problem source.}
\#\# Physical scenario

In figure, a block weighing \$14.0\textbackslash{} N\$, which can slide without friction on an incline at angle \$\textbackslash{}theta = 40.0\textasciicircum{}\textbackslash{}circ\$, is connected to the top of the incline by a massless spring of unstretched length \$0.450\textbackslash{} m\$ and spring constant \$120\textbackslash{} N/m\$.

\#\# Question

If the block is pulled slightly down the incline and released, what is the period of the resulting oscillations?

\#\# Answer choices

- A: A: 0.619 s

- B: B: 0.628 s

- C: C: 0.686 s

- D: D: 0.742 s

\#\# Figure caption (auxiliary; use the image as primary evidence)

The image depicts a physics diagram involving a spring, an inclined plane, and a block. Here are the details:

- **Inclined Plane**: The plane is angled, with an angle labeled as \textbackslash{}( \textbackslash{}theta \textbackslash{}).
- **Block**: A rectangular block rests on the inclined plane.
- **Spring**: A spring is attached to a fixed surface at the top of the incline and connected to the block.
- **Spring Constant**: The spring is labeled with a constant \textbackslash{}( k \textbackslash{}), representing the spring constant.
- **Relationships**: The block is in contact with the plane and the spring, suggesting a scenario involving elastic potential energy and gravitational forces.

This is a typical setup for analyzing mechanical systems involving friction, gravity, and spring force.

\#\# Dataset metadata

Wave Properties; Physical Model Grounding Reasoning, Multi-Formula Reasoning

\paragraph{Recorded answer/context.}
C: C: 0.686 s

\paragraph{Figure/image path.}
/root/proposal\_for\_physic/science-mango/phyx\_mini\_run/phyx\_data/test\_image/257.png

\paragraph{Formalization target.}
create a compiling Lean file with sorry bodies at `PhyXMiniProblems/problem\_phyx\_mini\_0257.lean`. The Lean declarations must preserve the physical quantities, dimensions or dimensional roles, figure labels, governing-law hypotheses, and final relation expressed by this problem.
Use Mathlib/Physlib names found through LeanExplore where available. If a physics API is missing, introduce faithful local abstractions rather than scalar placeholder aliases.

\begin{theorem}[Physics formalization target]
\label{thm:physics:phyx_mini_0257:target}
\lean{PhyXMiniProblems.ProblemPhyXMini0257.oscillationPeriod_matches_recordedAnswerC}
\uses{def:physics:phyx-mini-0257:phyxminiproblems-problemphyxmini0257-inclinedspringsetup, def:physics:phyx-mini-0257:phyxminiproblems-problemphyxmini0257-matchesproblemdata, def:physics:phyx-mini-0257:phyxminiproblems-problemphyxmini0257-matchessuppliedfigure, def:physics:phyx-mini-0257:phyxminiproblems-problemphyxmini0257-satisfiesinclinedspringoscillatorlaws, def:physics:phyx-mini-0257:phyxminiproblems-problemphyxmini0257-recordedanswerchoice, def:physics:phyx-mini-0257:phyxminiproblems-problemphyxmini0257-matchesdisplayedperiod, def:physics:phyx-mini-0257:phyxminiproblems-problemphyxmini0257-isuniqueclosestanswer}
The assigned autoformalize agent should translate this physics problem into Lean declarations in the covered file, with theorem and lemma proof bodies written as `by sorry`.
\end{theorem}
\begin{proof}
This is an autoformalization task, not a proof task. Produce statements that can later be proved by the physics prover without weakening the physical model.
\end{proof}

% --- Archon declaration topology begin ---
\section{Declaration topology}
\begin{definition}[Mass Quantity]
  \label{def:physics:phyx-mini-0257:phyxminiproblems-problemphyxmini0257-massquantity}
  \lean{PhyXMiniProblems.ProblemPhyXMini0257.MassQuantity}
  A physical mass
\end{definition}
\begin{proof}
  This is the declared model definition of Mass Quantity.
\end{proof}
\begin{definition}[Force Quantity]
  \label{def:physics:phyx-mini-0257:phyxminiproblems-problemphyxmini0257-forcequantity}
  \lean{PhyXMiniProblems.ProblemPhyXMini0257.ForceQuantity}
  A physical force, used for the block's weight
\end{definition}
\begin{proof}
  This is the declared model definition of Force Quantity.
\end{proof}
\begin{definition}[Length Quantity]
  \label{def:physics:phyx-mini-0257:phyxminiproblems-problemphyxmini0257-lengthquantity}
  \lean{PhyXMiniProblems.ProblemPhyXMini0257.LengthQuantity}
  A physical length
\end{definition}
\begin{proof}
  This is the declared model definition of Length Quantity.
\end{proof}
\begin{definition}[Acceleration Quantity]
  \label{def:physics:phyx-mini-0257:phyxminiproblems-problemphyxmini0257-accelerationquantity}
  \lean{PhyXMiniProblems.ProblemPhyXMini0257.AccelerationQuantity}
  A physical acceleration
\end{definition}
\begin{proof}
  This is the declared model definition of Acceleration Quantity.
\end{proof}
\begin{definition}[Spring Constant Quantity]
  \label{def:physics:phyx-mini-0257:phyxminiproblems-problemphyxmini0257-springconstantquantity}
  \lean{PhyXMiniProblems.ProblemPhyXMini0257.SpringConstantQuantity}
  A spring stiffness, with dimension force per length
\end{definition}
\begin{proof}
  This is the declared model definition of Spring Constant Quantity.
\end{proof}
\begin{definition}[Speed Quantity]
  \label{def:physics:phyx-mini-0257:phyxminiproblems-problemphyxmini0257-speedquantity}
  \lean{PhyXMiniProblems.ProblemPhyXMini0257.SpeedQuantity}
  A physical speed, used for the release speed
\end{definition}
\begin{proof}
  This is the declared model definition of Speed Quantity.
\end{proof}
\begin{definition}[Time Quantity]
  \label{def:physics:phyx-mini-0257:phyxminiproblems-problemphyxmini0257-timequantity}
  \lean{PhyXMiniProblems.ProblemPhyXMini0257.TimeQuantity}
  A physical duration
\end{definition}
\begin{proof}
  This is the declared model definition of Time Quantity.
\end{proof}
\begin{definition}[mass In Kilograms]
  \label{def:physics:phyx-mini-0257:phyxminiproblems-problemphyxmini0257-massinkilograms}
  \lean{PhyXMiniProblems.ProblemPhyXMini0257.massInKilograms}
  \uses{def:physics:phyx-mini-0257:phyxminiproblems-problemphyxmini0257-massquantity}
  Kilogram readout of a physical mass
\end{definition}
\begin{proof}
  This is the declared model definition of mass In Kilograms.
\end{proof}
\begin{definition}[force In Newtons]
  \label{def:physics:phyx-mini-0257:phyxminiproblems-problemphyxmini0257-forceinnewtons}
  \lean{PhyXMiniProblems.ProblemPhyXMini0257.forceInNewtons}
  \uses{def:physics:phyx-mini-0257:phyxminiproblems-problemphyxmini0257-forcequantity}
  Newton readout of a physical force
\end{definition}
\begin{proof}
  This is the declared model definition of force In Newtons.
\end{proof}
\begin{definition}[length In Meters]
  \label{def:physics:phyx-mini-0257:phyxminiproblems-problemphyxmini0257-lengthinmeters}
  \lean{PhyXMiniProblems.ProblemPhyXMini0257.lengthInMeters}
  \uses{def:physics:phyx-mini-0257:phyxminiproblems-problemphyxmini0257-lengthquantity}
  Metre readout of a physical length
\end{definition}
\begin{proof}
  This is the declared model definition of length In Meters.
\end{proof}
\begin{definition}[acceleration In Meters Per Second Squared]
  \label{def:physics:phyx-mini-0257:phyxminiproblems-problemphyxmini0257-accelerationinmeterspersecondsquared}
  \lean{PhyXMiniProblems.ProblemPhyXMini0257.accelerationInMetersPerSecondSquared}
  \uses{def:physics:phyx-mini-0257:phyxminiproblems-problemphyxmini0257-accelerationquantity}
  Metres-per-second-squared readout of a physical acceleration
\end{definition}
\begin{proof}
  This is the declared model definition of acceleration In Meters Per Second Squared.
\end{proof}
\begin{definition}[spring Constant In Newtons Per Meter]
  \label{def:physics:phyx-mini-0257:phyxminiproblems-problemphyxmini0257-springconstantinnewtonspermeter}
  \lean{PhyXMiniProblems.ProblemPhyXMini0257.springConstantInNewtonsPerMeter}
  \uses{def:physics:phyx-mini-0257:phyxminiproblems-problemphyxmini0257-springconstantquantity}
  Newton-per-metre readout of a spring stiffness
\end{definition}
\begin{proof}
  This is the declared model definition of spring Constant In Newtons Per Meter.
\end{proof}
\begin{definition}[speed In Meters Per Second]
  \label{def:physics:phyx-mini-0257:phyxminiproblems-problemphyxmini0257-speedinmeterspersecond}
  \lean{PhyXMiniProblems.ProblemPhyXMini0257.speedInMetersPerSecond}
  \uses{def:physics:phyx-mini-0257:phyxminiproblems-problemphyxmini0257-speedquantity}
  Metres-per-second readout of a physical speed
\end{definition}
\begin{proof}
  This is the declared model definition of speed In Meters Per Second.
\end{proof}
\begin{definition}[time In Seconds]
  \label{def:physics:phyx-mini-0257:phyxminiproblems-problemphyxmini0257-timeinseconds}
  \lean{PhyXMiniProblems.ProblemPhyXMini0257.timeInSeconds}
  \uses{def:physics:phyx-mini-0257:phyxminiproblems-problemphyxmini0257-timequantity}
  Second readout of a physical duration
\end{definition}
\begin{proof}
  This is the declared model definition of time In Seconds.
\end{proof}
\begin{definition}[degrees To Radians]
  \label{def:physics:phyx-mini-0257:phyxminiproblems-problemphyxmini0257-degreestoradians}
  \lean{PhyXMiniProblems.ProblemPhyXMini0257.degreesToRadians}
  Convert a dimensionless angle stated in degrees to radians
\end{definition}
\begin{proof}
  This is the declared model definition of degrees To Radians.
\end{proof}
\begin{definition}[Block Contact Model]
  \label{def:physics:phyx-mini-0257:phyxminiproblems-problemphyxmini0257-blockcontactmodel}
  \lean{PhyXMiniProblems.ProblemPhyXMini0257.BlockContactModel}
  Contact model between the block and the incline
\end{definition}
\begin{proof}
  This is the declared model definition of Block Contact Model.
\end{proof}
\begin{definition}[Spring Mass Model]
  \label{def:physics:phyx-mini-0257:phyxminiproblems-problemphyxmini0257-springmassmodel}
  \lean{PhyXMiniProblems.ProblemPhyXMini0257.SpringMassModel}
  Mass model for the connecting spring
\end{definition}
\begin{proof}
  This is the declared model definition of Spring Mass Model.
\end{proof}
\begin{definition}[Incline Direction]
  \label{def:physics:phyx-mini-0257:phyxminiproblems-problemphyxmini0257-inclinedirection}
  \lean{PhyXMiniProblems.ProblemPhyXMini0257.InclineDirection}
  Direction in which the block is initially pulled
\end{definition}
\begin{proof}
  This is the declared model definition of Incline Direction.
\end{proof}
\begin{definition}[Displacement Regime]
  \label{def:physics:phyx-mini-0257:phyxminiproblems-problemphyxmini0257-displacementregime}
  \lean{PhyXMiniProblems.ProblemPhyXMini0257.DisplacementRegime}
  Qualitative displacement regime expressed by the word "slightly"
\end{definition}
\begin{proof}
  This is the declared model definition of Displacement Regime.
\end{proof}
\begin{definition}[Figure Component]
  \label{def:physics:phyx-mini-0257:phyxminiproblems-problemphyxmini0257-figurecomponent}
  \lean{PhyXMiniProblems.ProblemPhyXMini0257.FigureComponent}
  Components visible in the supplied diagram
\end{definition}
\begin{proof}
  This is the declared model definition of Figure Component.
\end{proof}
\begin{definition}[Figure Label]
  \label{def:physics:phyx-mini-0257:phyxminiproblems-problemphyxmini0257-figurelabel}
  \lean{PhyXMiniProblems.ProblemPhyXMini0257.FigureLabel}
  Symbols printed in the supplied diagram
\end{definition}
\begin{proof}
  This is the declared model definition of Figure Label.
\end{proof}
\begin{definition}[Spring Orientation]
  \label{def:physics:phyx-mini-0257:phyxminiproblems-problemphyxmini0257-springorientation}
  \lean{PhyXMiniProblems.ProblemPhyXMini0257.SpringOrientation}
  Orientation of the spring relative to the incline
\end{definition}
\begin{proof}
  This is the declared model definition of Spring Orientation.
\end{proof}
\begin{definition}[Spring Upper Attachment]
  \label{def:physics:phyx-mini-0257:phyxminiproblems-problemphyxmini0257-springupperattachment}
  \lean{PhyXMiniProblems.ProblemPhyXMini0257.SpringUpperAttachment}
  Upper attachment shown for the spring
\end{definition}
\begin{proof}
  This is the declared model definition of Spring Upper Attachment.
\end{proof}
\begin{definition}[Spring Lower Attachment]
  \label{def:physics:phyx-mini-0257:phyxminiproblems-problemphyxmini0257-springlowerattachment}
  \lean{PhyXMiniProblems.ProblemPhyXMini0257.SpringLowerAttachment}
  Lower attachment shown for the spring
\end{definition}
\begin{proof}
  This is the declared model definition of Spring Lower Attachment.
\end{proof}
\begin{definition}[Supplied Figure]
  \label{def:physics:phyx-mini-0257:phyxminiproblems-problemphyxmini0257-suppliedfigure}
  \lean{PhyXMiniProblems.ProblemPhyXMini0257.SuppliedFigure}
  \uses{def:physics:phyx-mini-0257:phyxminiproblems-problemphyxmini0257-figurecomponent, def:physics:phyx-mini-0257:phyxminiproblems-problemphyxmini0257-figurelabel, def:physics:phyx-mini-0257:phyxminiproblems-problemphyxmini0257-springorientation, def:physics:phyx-mini-0257:phyxminiproblems-problemphyxmini0257-springupperattachment, def:physics:phyx-mini-0257:phyxminiproblems-problemphyxmini0257-springlowerattachment}
  Structured readout of the primary image
\end{definition}
\begin{proof}
  This is the declared model definition of Supplied Figure.
\end{proof}
\begin{definition}[Inclined Spring Setup]
  \label{def:physics:phyx-mini-0257:phyxminiproblems-problemphyxmini0257-inclinedspringsetup}
  \lean{PhyXMiniProblems.ProblemPhyXMini0257.InclinedSpringSetup}
  \uses{def:physics:phyx-mini-0257:phyxminiproblems-problemphyxmini0257-massquantity, def:physics:phyx-mini-0257:phyxminiproblems-problemphyxmini0257-forcequantity, def:physics:phyx-mini-0257:phyxminiproblems-problemphyxmini0257-lengthquantity, def:physics:phyx-mini-0257:phyxminiproblems-problemphyxmini0257-accelerationquantity, def:physics:phyx-mini-0257:phyxminiproblems-problemphyxmini0257-springconstantquantity, def:physics:phyx-mini-0257:phyxminiproblems-problemphyxmini0257-speedquantity, def:physics:phyx-mini-0257:phyxminiproblems-problemphyxmini0257-timequantity, def:physics:phyx-mini-0257:phyxminiproblems-problemphyxmini0257-blockcontactmodel, def:physics:phyx-mini-0257:phyxminiproblems-problemphyxmini0257-springmassmodel, def:physics:phyx-mini-0257:phyxminiproblems-problemphyxmini0257-inclinedirection, def:physics:phyx-mini-0257:phyxminiproblems-problemphyxmini0257-displacementregime, def:physics:phyx-mini-0257:phyxminiproblems-problemphyxmini0257-suppliedfigure}
  The physical system and its unknown oscillation period. The physical period is an independent field: it is neither defined to be a displayed answer nor assigned a numerical value in this structure
\end{definition}
\begin{proof}
  This is the declared model definition of Inclined Spring Setup.
\end{proof}
\begin{definition}[Matches Problem Data]
  \label{def:physics:phyx-mini-0257:phyxminiproblems-problemphyxmini0257-matchesproblemdata}
  \lean{PhyXMiniProblems.ProblemPhyXMini0257.MatchesProblemData}
  \uses{def:physics:phyx-mini-0257:phyxminiproblems-problemphyxmini0257-forceinnewtons, def:physics:phyx-mini-0257:phyxminiproblems-problemphyxmini0257-lengthinmeters, def:physics:phyx-mini-0257:phyxminiproblems-problemphyxmini0257-accelerationinmeterspersecondsquared, def:physics:phyx-mini-0257:phyxminiproblems-problemphyxmini0257-springconstantinnewtonspermeter, def:physics:phyx-mini-0257:phyxminiproblems-problemphyxmini0257-speedinmeterspersecond, def:physics:phyx-mini-0257:phyxminiproblems-problemphyxmini0257-inclinedspringsetup}
  Numerical and qualitative data stated in the problem. The value 9.8 m/s² is the standard gravitational calibration needed to infer mass from the given weight. No period or answer-choice value occurs here
\end{definition}
\begin{proof}
  This is the declared model definition of Matches Problem Data.
\end{proof}
\begin{definition}[Matches Supplied Figure]
  \label{def:physics:phyx-mini-0257:phyxminiproblems-problemphyxmini0257-matchessuppliedfigure}
  \lean{PhyXMiniProblems.ProblemPhyXMini0257.MatchesSuppliedFigure}
  \uses{def:physics:phyx-mini-0257:phyxminiproblems-problemphyxmini0257-inclinedspringsetup}
  Primary-image readout: an inclined plane, block, spring, and top anchor are visible; θ labels the incline and k labels the spring; the spring runs along the incline from the block to a fixed upper anchor
\end{definition}
\begin{proof}
  This is the declared model definition of Matches Supplied Figure.
\end{proof}
\begin{definition}[Satisfies Inclined Spring Oscillator Laws]
  \label{def:physics:phyx-mini-0257:phyxminiproblems-problemphyxmini0257-satisfiesinclinedspringoscillatorlaws}
  \lean{PhyXMiniProblems.ProblemPhyXMini0257.SatisfiesInclinedSpringOscillatorLaws}
  \uses{def:physics:phyx-mini-0257:phyxminiproblems-problemphyxmini0257-massinkilograms, def:physics:phyx-mini-0257:phyxminiproblems-problemphyxmini0257-forceinnewtons, def:physics:phyx-mini-0257:phyxminiproblems-problemphyxmini0257-lengthinmeters, def:physics:phyx-mini-0257:phyxminiproblems-problemphyxmini0257-accelerationinmeterspersecondsquared, def:physics:phyx-mini-0257:phyxminiproblems-problemphyxmini0257-springconstantinnewtonspermeter, def:physics:phyx-mini-0257:phyxminiproblems-problemphyxmini0257-timeinseconds, def:physics:phyx-mini-0257:phyxminiproblems-problemphyxmini0257-degreestoradians, def:physics:phyx-mini-0257:phyxminiproblems-problemphyxmini0257-inclinedspringsetup}
  The governing laws for the frictionless, small-displacement motion. The first two equations identify weight with m g and locate the statically shifted equilibrium using the component of gravity along the incline. The last three fields connect the physical system to Physlib's positive-mass, positive-stiffness harmonic oscillator and its general period. None assumes the requested numerical period or a displayed answer
\end{definition}
\begin{proof}
  This is the declared model definition of Satisfies Inclined Spring Oscillator Laws.
\end{proof}
\begin{lemma}[oscillation Period formula]
  \label{lem:physics:phyx-mini-0257:phyxminiproblems-problemphyxmini0257-oscillationperiod-formula}
  \lean{PhyXMiniProblems.ProblemPhyXMini0257.oscillationPeriod_formula}
  \uses{def:physics:phyx-mini-0257:phyxminiproblems-problemphyxmini0257-massinkilograms, def:physics:phyx-mini-0257:phyxminiproblems-problemphyxmini0257-springconstantinnewtonspermeter, def:physics:phyx-mini-0257:phyxminiproblems-problemphyxmini0257-timeinseconds, def:physics:phyx-mini-0257:phyxminiproblems-problemphyxmini0257-inclinedspringsetup, def:physics:phyx-mini-0257:phyxminiproblems-problemphyxmini0257-satisfiesinclinedspringoscillatorlaws}
  The Physlib oscillator laws give the standard mass--spring period formula. Gravity and incline angle affect the equilibrium extension but not this linearized period
\end{lemma}
\begin{proof}
  The conclusion follows from the governing relations and figure readouts named in the statement of oscillation Period formula.
\end{proof}
\begin{definition}[Answer Choice]
  \label{def:physics:phyx-mini-0257:phyxminiproblems-problemphyxmini0257-answerchoice}
  \lean{PhyXMiniProblems.ProblemPhyXMini0257.AnswerChoice}
  Labels of the four displayed answers
\end{definition}
\begin{proof}
  This is the declared model definition of Answer Choice.
\end{proof}
\begin{definition}[seconds]
  \label{def:physics:phyx-mini-0257:phyxminiproblems-problemphyxmini0257-answerchoice-seconds}
  \lean{PhyXMiniProblems.ProblemPhyXMini0257.AnswerChoice.seconds}
  \uses{def:physics:phyx-mini-0257:phyxminiproblems-problemphyxmini0257-answerchoice}
  Period in seconds printed beside each answer label
\end{definition}
\begin{proof}
  This is the declared model definition of seconds.
\end{proof}
\begin{definition}[recorded Answer Choice]
  \label{def:physics:phyx-mini-0257:phyxminiproblems-problemphyxmini0257-recordedanswerchoice}
  \lean{PhyXMiniProblems.ProblemPhyXMini0257.recordedAnswerChoice}
  \uses{def:physics:phyx-mini-0257:phyxminiproblems-problemphyxmini0257-answerchoice}
  The answer label recorded in the supplied dataset metadata
\end{definition}
\begin{proof}
  This is the declared model definition of recorded Answer Choice.
\end{proof}
\begin{definition}[Matches Displayed Period]
  \label{def:physics:phyx-mini-0257:phyxminiproblems-problemphyxmini0257-matchesdisplayedperiod}
  \lean{PhyXMiniProblems.ProblemPhyXMini0257.MatchesDisplayedPeriod}
  \uses{def:physics:phyx-mini-0257:phyxminiproblems-problemphyxmini0257-timequantity, def:physics:phyx-mini-0257:phyxminiproblems-problemphyxmini0257-timeinseconds, def:physics:phyx-mini-0257:phyxminiproblems-problemphyxmini0257-answerchoice}
  Agreement with a displayed period rounded to the nearest millisecond
\end{definition}
\begin{proof}
  This is the declared model definition of Matches Displayed Period.
\end{proof}
\begin{definition}[Is Unique Closest Answer]
  \label{def:physics:phyx-mini-0257:phyxminiproblems-problemphyxmini0257-isuniqueclosestanswer}
  \lean{PhyXMiniProblems.ProblemPhyXMini0257.IsUniqueClosestAnswer}
  \uses{def:physics:phyx-mini-0257:phyxminiproblems-problemphyxmini0257-timequantity, def:physics:phyx-mini-0257:phyxminiproblems-problemphyxmini0257-timeinseconds, def:physics:phyx-mini-0257:phyxminiproblems-problemphyxmini0257-answerchoice}
  The selected choice is strictly closer than every other displayed value
\end{definition}
\begin{proof}
  This is the declared model definition of Is Unique Closest Answer.
\end{proof}
% --- Archon declaration topology end ---
% --- Archon physics formalization source end ---
